\documentclass[11pt]{article}
% Compile with pdfLaTeX from the repository root or from this folder.
\makeatletter
\def\input@path{{../}}
\makeatother
\usepackage{notes}
\setboolean{showanswer}{true}
\renewcommand{\solutionheight}{22mm}

\begin{document}
\lesson[Stewart, Calculus, 9th edition\enspace\textbullet\enspace Section 12.4]
  {1.4}{Cross Product}

\section{A vector perpendicular to two vectors}
The dot product measures alignment and returns a scalar. The cross product
of two vectors in $\mathbb{R}^3$ returns a \textbf{vector} perpendicular to both.

\begin{definition}[Cross product]
For $\B{u}=\cvec{u_1,u_2,u_3}$ and $\B{v}=\cvec{v_1,v_2,v_3}$,
\[
 \B{u}\times\B{v}
 =\cvec{u_2v_3-u_3v_2,\ u_3v_1-u_1v_3,\ u_1v_2-u_2v_1}.
\]
The order matters: $\B{v}\times\B{u}=-(\B{u}\times\B{v})$.
\end{definition}

\subsection*{A determinant helps organize the calculation}
We use the formal determinant notation
\[
 \B{u}\times\B{v}
 =\begin{vmatrix}
 \B{i}&\B{j}&\B{k}\\
 u_1&u_2&u_3\\
 v_1&v_2&v_3
 \end{vmatrix}.
\]
Expand along the first row with signs $+,-,+$. The minus sign on the
$\B{j}$ term is easy to miss:
\[
 \B{u}\times\B{v}
 =(u_2v_3-u_3v_2)\B{i}
 -(u_1v_3-u_3v_1)\B{j}
 +(u_1v_2-u_2v_1)\B{k}.
\]

\begin{example}{Find $\B{u}\times\B{v}$ for
$\B{u}=\cvec{1,2,0}$ and $\B{v}=\cvec{0,1,2}$.}
Expanding the determinant gives
\begin{align*}
 \B{u}\times\B{v}
 &=\begin{vmatrix}\B{i}&\B{j}&\B{k}\\1&2&0\\0&1&2\end{vmatrix}\\
 &=(4-0)\B{i}-(2-0)\B{j}+(1-0)\B{k}=\cvec{4,-2,1}.
\end{align*}
Check perpendicularity using dot products:
\[
 \cvec{4,-2,1}\cdot\cvec{1,2,0}=4-4=0,\qquad
 \cvec{4,-2,1}\cdot\cvec{0,1,2}=-2+2=0.
\]
\end{example}

\subsection*{Useful properties}
For vectors $\B{u},\B{v},\B{w}$ and a scalar $c$,
\[
 \B{u}\times(\B{v}+\B{w})=\B{u}\times\B{v}+\B{u}\times\B{w},
 \qquad
 (c\B{u})\times\B{v}=c(\B{u}\times\B{v}).
\]
Also, $\B{u}\times\B{u}=\B{0}$. Unlike ordinary multiplication, the cross
product is neither commutative nor associative.

\section{Direction, length, and area}
\begin{definition}[Geometric meaning]
For nonzero vectors $\B{u}$ and $\B{v}$ with angle $0\leq\theta\leq\pi$,
\[
 \lVert\B{u}\times\B{v}\rVert
 =\lVert\B{u}\rVert\lVert\B{v}\rVert\sin\theta.
\]
When the cross product is nonzero, its direction is given by the
\textbf{right-hand rule}: curl your fingers from $\B{u}$ toward $\B{v}$
through the smaller angle; your thumb points along $\B{u}\times\B{v}$.
\end{definition}

\begin{center}
\begin{tikzpicture}[x={(-0.75cm,-0.4cm)},y={(0.9cm,-0.12cm)},z={(0cm,0.85cm)},
 >=Latex,font=\small]
 \fill[notesblue!5] (0,0,0)--(2.8,0,0)--(2.8,3,0)--(0,3,0)--cycle;
 \draw[dashed,notesrule] (2.8,0,0)--(2.8,3,0)--(0,3,0);
 \draw[->,very thick,notesblue] (0,0,0)--(2.8,0,0) node[below left]{$\B{u}$};
 \draw[->,very thick,notesteal] (0,0,0)--(0,3,0) node[right]{$\B{v}$};
 \draw[->,very thick,notesink] (0,0,0)--(0,0,2.1) node[above]{$\B{u}\times\B{v}$};
 \fill[notesink] (0,0,0) circle (1.6pt);
\end{tikzpicture}
\end{center}

In particular, $\B{i}\times\B{j}=\B{k}$,
$\B{j}\times\B{k}=\B{i}$, and $\B{k}\times\B{i}=\B{j}$.
Reversing any pair changes the sign. For two nonzero vectors, the cross
product is zero exactly when they are parallel, including opposite directions.

\begin{example}{Find a unit vector perpendicular to both
$\B{u}=\cvec{1,2,0}$ and $\B{v}=\cvec{0,1,2}$.}
From Example 1, a perpendicular vector is $\B{n}=\cvec{4,-2,1}$.
Its magnitude is $\sqrt{16+4+1}=\sqrt{21}$, so one unit normal is
\[
 \widehat{\B{n}}=\frac{1}{\sqrt{21}}\cvec{4,-2,1}.
\]
The opposite vector $-\widehat{\B{n}}$ is also a unit normal.
\end{example}

\subsection*{Areas from two sides}
A parallelogram with adjacent sides $\B{u}$ and $\B{v}$ has base
$\lVert\B{u}\rVert$ and height $\lVert\B{v}\rVert\sin\theta$. Therefore,
\[
 A_{\text{parallelogram}}=\lVert\B{u}\times\B{v}\rVert,
 \qquad A_{\text{triangle}}=\tfrac12\lVert\B{u}\times\B{v}\rVert.
\]

\begin{example}{Find the area of the triangle with vertices
$P=(1,0,0)$, $Q=(2,2,0)$, and $R=(1,1,2)$.}
Use two sides starting at the same vertex:
$\vec{PQ}=\cvec{1,2,0}$ and $\vec{PR}=\cvec{0,1,2}$. Thus
\[
 A=\tfrac12\lVert\vec{PQ}\times\vec{PR}\rVert=\frac{\sqrt{21}}2.
\]
\end{example}

\newpage
\section{Torque}
A force can turn an object about a pivot. Both the force and where it is
applied affect this turning effect.

\begin{definition}[Torque]
Let $\B{r}$ point from the pivot to the point where a force $\B{F}$ is applied.
The \textbf{torque} about the pivot is
\[
 \boldsymbol{\tau}=\B{r}\times\B{F},\qquad
 \lVert\boldsymbol{\tau}\rVert
 =\lVert\B{r}\rVert\lVert\B{F}\rVert\sin\theta.
\]
Its direction gives the axis of rotation by the right-hand rule.
\end{definition}

\begin{example}{A force $\B{F}=\cvec{0,3,0}$ newtons is applied at position
$\B{r}=\cvec{2,0,0}$ meters relative to a pivot. Find the torque.}
Since $\B{r}=2\B{i}$ and $\B{F}=3\B{j}$,
\[
 \boldsymbol{\tau}=(2\B{i})\times(3\B{j})=6\B{k}.
\]
The torque has magnitude $6$ newton-meters and points along the positive
$z$-axis. A force parallel to $\B{r}$ would produce zero torque.
\end{example}

\section{The scalar triple product}
Combining a cross product with a dot product gives a volume calculation.

\begin{definition}[Scalar triple product]
The \textbf{scalar triple product} of $\B{a}$, $\B{b}$, and $\B{c}$ is
\[
 \B{a}\cdot(\B{b}\times\B{c}).
\]
The volume of the parallelepiped with these three adjacent edge vectors is
\[
 V=\bigl|\B{a}\cdot(\B{b}\times\B{c})\bigr|.
\]
\end{definition}

The cross product $\B{b}\times\B{c}$ supplies a base area and a normal
direction. Dotting with $\B{a}$ accounts for the perpendicular height.
The absolute value makes the volume nonnegative.

\begin{example}{Find the volume determined by
$\B{a}=\cvec{1,0,1}$, $\B{b}=\cvec{0,2,0}$, and $\B{c}=\cvec{1,0,3}$.}
First compute
\[
 \B{b}\times\B{c}
 =\begin{vmatrix}\B{i}&\B{j}&\B{k}\\0&2&0\\1&0&3\end{vmatrix}
 =\cvec{6,0,-2}.
\]
Then $V=|\cvec{1,0,1}\cdot\cvec{6,0,-2}|=|6-2|=4$ cubic units.
\end{example}

Three vectors are \textbf{coplanar} when they lie in a common plane through
the origin. This happens exactly when their scalar triple product is zero:
the corresponding solid has zero volume.
\end{document}
